% =============================================================================
% Coxeter diagrams: constructors
% =============================================================================
% Building blocks that draw a Coxeter diagram from its data: roots (vertices)
% with their marks and positions, edges with their Coxeter weights, chains,
% cycles and maximal parabolic subdiagrams. The diagrams themselves are
% objects, one file each, under figures/objects/coxeter/, assembled from these
% constructors; a figure places one with
%   \pic (A) at (6,0) [<keys>] {object=coxeter/sterk-cusp-1};
% and draws on its roots A-0, A-1, ...:
%   \scoped[on background layer] \draw[parabolic] (A-0.center) -- (A-8.center);
% \input by dzg-constructors.tex, so "@" is a letter here.
%
% Commands (the root <i> is the node -<i>, labelled by the label scheme):
%   \coxeterlabels{<scheme>}
%       the object's canonical label scheme, used unless the figure passes
%       `root labels`; an object states it before its first root.
%   \coxeterroot{<i>}{<mark>}{<coordinate>}{<label direction>}
%       one root. <mark> is a `vertex` style (white, black, doubled) or even.
%   \coxeterroots{<i>/<mark>/<coordinate>/<label direction>, ...}
%       several roots. Brace a coordinate that contains a comma; place a calc
%       coordinate with \coxeterroot, since \foreach keeps the braces.
%   \coxeteredges{<Coxeter weight>}{<path through -<i> names>}
%       edges of one weight (a `coxeter edge` value), on the edge layer.
%   \coxeterchain{<Coxeter weight>}{<first>}{<last>}
%       the edges i -- i+1 for first <= i < last.
%   \coxetersquare{<side>}{<edges per side>}{<first>}{<sw|nw>}
%       {<mark at corners and every second root>}{<mark between>}
%       a cycle of simple edges around the square [0,side]^2, numbered from
%       <first> at the south-west corner counterclockwise (sw) or at the
%       north-west corner clockwise (nw), labelled outward.
%   \coxeterEdiagram{<n>}{<mark>}   (dzg-vinberg.tex)
%   \coxeterparabolic{<name>}{{<i>,<j>,...}, {...}, ...}
%       a maximal parabolic subdiagram of the object: the highlight stroked
%       through the root centres of each chain, drawn when the figure passes
%       `maximal parabolic=<name>`. The vertex sets are data of the object.
%   \coxeterpair[<keys>], \coxetervertex{<mark>}   table legends
%
% Pics:
%   dynkin folding={<unfolded options>}{<diagram>}{<folded options>}{<diagram>}
%   root chain={<mark>/<norm label>/<weight to the next root>, ...}
%
% Keys:
%   root labels=alpha|r|ell|index|none   the label scheme (see below); a bare
%                                        `root labels` selects the canonical one
%   parity marks=true|false              even roots doubled, or white
%   maximal parabolic=<name>             highlight that parabolic subdiagram
%   vertex=even                          an even boundary root (AEGS25
%                                        Sec. 3.5): doubled under `parity
%                                        marks`, white otherwise
%
% Label schemes; the index is always the canonical root name:
%   alpha  $\alpha_i$
%   r      $r_i$
%   ell    $\ell_i$    the coordinate \ell_i = \lambda\cdot\alpha_i of an
%                      integral-affine sphere B(\lambda) (IAS figures)
%   index  $i$         the bare index
%   none   no labels

\newif\ifdzg@rootlabels
\newif\ifdzg@paritymarks
\dzg@paritymarkstrue
\def\dzg@unset{unset}
\def\dzg@labelscheme{unset}
\def\dzg@parabolic{}
\tikzset{
  root labels/.is choice,
  root labels/canonical/.code={\dzg@rootlabelstrue\def\dzg@labelscheme{unset}},
  root labels/alpha/.code={\dzg@rootlabelstrue\def\dzg@labelscheme{alpha}},
  root labels/r/.code={\dzg@rootlabelstrue\def\dzg@labelscheme{r}},
  root labels/ell/.code={\dzg@rootlabelstrue\def\dzg@labelscheme{ell}},
  root labels/index/.code={\dzg@rootlabelstrue\def\dzg@labelscheme{index}},
  root labels/none/.code={\dzg@rootlabelsfalse\def\dzg@labelscheme{none}},
  root labels/.default=canonical,
  parity marks/.is if=dzg@paritymarks,
  vertex/even/.code={%
    \ifdzg@paritymarks\pgfkeysalso{vertex/doubled}\else\pgfkeysalso{vertex/white}\fi},
  maximal parabolic/.code={\def\dzg@parabolic{#1}},
}
\def\dzg@label@alpha#1{$\alpha_{#1}$}
\def\dzg@label@r#1{$r_{#1}$}
\def\dzg@label@ell#1{$\ell_{#1}$}
\def\dzg@label@index#1{$#1$}
\def\dzg@label@unset#1{$r_{#1}$}
\def\dzg@rootlabel#1{\csname dzg@label@\dzg@labelscheme\endcsname{#1}}

\newcommand\coxeterlabels[1]{%
  \ifx\dzg@labelscheme\dzg@unset\dzg@rootlabelstrue\def\dzg@labelscheme{#1}\fi}

\newcommand\coxeterroot[4]{%
  \ifdzg@rootlabels
    \node[vertex=#2, label={[vertex label]#4:\dzg@rootlabel{#1}}] (-#1) at (#3) {};%
  \else
    \node[vertex=#2] (-#1) at (#3) {};%
  \fi}

\newcommand\coxeterroots[1]{%
  \foreach \dzg@i/\dzg@m/\dzg@p/\dzg@a in {#1}{\coxeterroot{\dzg@i}{\dzg@m}{\dzg@p}{\dzg@a}}}

\newcommand\coxeteredges[2]{\scoped[on edge layer] \draw[coxeter edge=#1] #2;}

% A path \foreach through node names restarts every segment at the first node,
% so each edge is its own path.
\newcommand\coxeterchain[3]{%
  \pgfmathtruncatemacro\dzg@stop{#3-1}%
  \foreach \dzg@k in {#2,...,\dzg@stop}{%
    \pgfmathtruncatemacro\dzg@n{\dzg@k+1}%
    \coxeteredges{#1}{(-\dzg@k) -- (-\dzg@n)}}}

\newcommand\coxetersquare[6]{%
  \def\dzg@L{#1}%
  \pgfmathtruncatemacro\dzg@last{#3+4*#2-1}%
  \foreach \dzg@k in {#3,...,\dzg@last}{%
    \pgfmathtruncatemacro\dzg@s{div(\dzg@k-#3,#2)}%
    \pgfmathtruncatemacro\dzg@t{mod(\dzg@k-#3,#2)}%
    \pgfmathsetmacro\dzg@u{#1*\dzg@t/#2}%
    \pgfmathsetmacro\dzg@v{#1-\dzg@u}%
    \csname dzg@square@#4\endcsname
    \ifodd\dzg@t\def\dzg@m{#6}\else\def\dzg@m{#5}\fi
    \coxeterroot{\dzg@k}{\dzg@m}{\dzg@p}{\dzg@a}}%
  \coxeterchain{3}{#3}{\dzg@last}%
  \coxeteredges{3}{(-\dzg@last) -- (-#3)}}
% Position \dzg@p and outward label direction \dzg@a of the root at step
% \dzg@t along side \dzg@s.
\def\dzg@square@corner#1#2{\ifnum\dzg@t=0 \def\dzg@a{#1}\else\def\dzg@a{#2}\fi}
\def\dzg@square@sw{%
  \ifcase\dzg@s
    \edef\dzg@p{\dzg@u,0}\dzg@square@corner{225}{270}%
  \or\edef\dzg@p{\dzg@L,\dzg@u}\dzg@square@corner{315}{0}%
  \or\edef\dzg@p{\dzg@v,\dzg@L}\dzg@square@corner{45}{90}%
  \or\edef\dzg@p{0,\dzg@v}\dzg@square@corner{135}{180}%
  \fi}
\def\dzg@square@nw{%
  \ifcase\dzg@s
    \edef\dzg@p{\dzg@u,\dzg@L}\dzg@square@corner{135}{90}%
  \or\edef\dzg@p{\dzg@L,\dzg@v}\dzg@square@corner{45}{0}%
  \or\edef\dzg@p{\dzg@v,0}\dzg@square@corner{315}{270}%
  \or\edef\dzg@p{0,\dzg@u}\dzg@square@corner{225}{180}%
  \fi}

% Each chain is its own stroke, starting with a zero-length segment so that a
% chain of one root is a dot.
\newcommand\coxeterparabolic[2]{%
  \def\dzg@name{#1}%
  \ifx\dzg@name\dzg@parabolic
    \scoped[on background layer]{%
      \foreach \dzg@c in {#2} {%
        \foreach \dzg@v [count=\dzg@k] in \dzg@c {%
          \ifnum\dzg@k=1 \xdef\dzg@path{(-\dzg@v.center) -- (-\dzg@v.center)}%
          \else\xdef\dzg@path{\dzg@path -- (-\dzg@v.center)}\fi}%
        \draw[parabolic] \dzg@path;}}%
  \fi}

% -----------------------------------------------------------------------------
% Folding of Dynkin diagrams (dynkin-diagrams package): the diagram with its
% symmetry above the folded diagram, joined by the quotient map.
%   \pic {dynkin folding={<unfolded options>}{<type><spec>}{<folded options>}
%     {<type><spec>}};
% e.g. {ply=2}{A{ooo.ooooooo.ooo}}{}{C{ooo.ooo*}}. The nodes -unfolded and
% -folded are the two diagrams.
% -----------------------------------------------------------------------------
% The folded diagram is left-aligned under the unfolded one. Its node resets
% `anchor`, which the nested picture of \dynkin would otherwise inherit.
\tikzset{pics/dynkin folding/.style n args={4}{code={%
  \node[inner sep=2pt] (-unfolded) at (0,0) {\dynkin[scale=3,#1]#2};
  \node[inner sep=2pt, anchor=north west] (-folded)
    at ($(-unfolded.south west)+(0,-10mm)$)
    {\tikzset{anchor=center}\dynkin[scale=3, arrows=false, #3]#4};
  \draw[quotient map] (-unfolded.south) -- (-unfolded.south |- -folded.north);}}}

% -----------------------------------------------------------------------------
% A linear Coxeter diagram with the norm of each root written above it:
%   \pic {root chain={<mark>/<norm label>/<weight to the next root>, ...,
%     <mark>/<norm label>/}};
% Weight 3 is a simple edge; the vertices are named -0, -1, ....
% -----------------------------------------------------------------------------
\tikzset{pics/root chain/.style={code={%
  \foreach \dzg@m/\dzg@l/\dzg@w [count=\dzg@k from 0] in {#1}{%
    \node[vertex=\dzg@m, label={[vertex label]90:{$\dzg@l$}}] (-\dzg@k) at (1.5*\dzg@k,0) {};
    \ifx\dzg@w\pgfutil@empty\else
      \coordinate (-\dzg@k-next) at (1.5*\dzg@k+1.5,0);
      \scoped[on edge layer] \draw[coxeter edge=\dzg@w] (-\dzg@k.center) -- (-\dzg@k-next);
    \fi}}}}

% -----------------------------------------------------------------------------
% Table legends: two roots and the edge between them, or one vertex.
% -----------------------------------------------------------------------------
%   \coxeterpair[marks=black/white, edge=4, roots=f_1/f_2, weight=m]
% marks: the two vertex marks (default white/white); edge: a coxeter edge
% value or none (default 3); roots: the two labels in math mode (default
% H_1/H_2); weight: text set above the edge (default empty).
%   \coxetervertex{white|black|doubled}
\pgfqkeys{/dzg/coxeter pair}{
  marks/.store in=\dzg@pairmarks, marks=white/white,
  edge/.store in=\dzg@pairedge, edge=3,
  roots/.store in=\dzg@pairroots, roots=H_1/H_2,
  weight/.store in=\dzg@pairweight, weight=,
}
\def\dzg@none{none}
\def\dzg@splitpair#1/#2\@nil#3#4{\def#3{#1}\def#4{#2}}
\newcommand{\coxeterpair}[1][]{%
  \begin{tikzpicture}[baseline=-0.5ex]
    \pgfqkeys{/dzg/coxeter pair}{#1}%
    \expandafter\dzg@splitpair\dzg@pairmarks\@nil\dzg@marka\dzg@markb
    \expandafter\dzg@splitpair\dzg@pairroots\@nil\dzg@roota\dzg@rootb
    \node[vertex=\dzg@marka, label={[vertex label]above:$\dzg@roota$}] (h1) at (0,0) {};
    \node[vertex=\dzg@markb, label={[vertex label]above:$\dzg@rootb$}] (h2) at (3,0) {};
    \ifx\dzg@pairedge\dzg@none\else
      \scoped[on edge layer] \draw[coxeter edge=\dzg@pairedge]
        (h1) -- node[vertex label, midway, above] {\dzg@pairweight} (h2);
    \fi
  \end{tikzpicture}}
\newcommand{\coxetervertex}[1]{\tikz[baseline=-0.5ex]\node[vertex=#1] {};}
